\tocchapter{Introduction}

\par Over the last two years, AI assistants have moved from answering questions to performing tasks. A coding assistant now opens files, runs tests and creates pull requests; an office assistant reads the mailbox and books meetings. To do this, the model needs a way to reach external systems. Until late 2024 every vendor built its own connectors. In November 2024 Anthropic published the \acrfull{glo:mcp}, an open specification that describes how an AI application finds the tools offered by an external program and how it calls them \cite{anthropic2024mcp}. The protocol was adopted quickly. OpenAI and Google announced support in the spring of 2025, and in December 2025 the project was transferred to the Agentic AI Foundation, which operates under the Linux Foundation \cite{aaif2025}. In July 2026 Netskope reported that \acrshort{glo:mcp} traffic among its customers had grown four times in only ten weeks \cite{netskope2026}.

\par The security of the ecosystem developed more slowly. The first version of the specification did not include any authorisation mechanism. Critical vulnerabilities were found in official reference servers and in the developer tools published by the maintainers, and internet-wide scans in 2026 found thousands of servers that accepted requests without checking who sent them \cite{zhou2026, padilla2026}. In May 2026 the U.S. National Security Agency released a separate information sheet on the secure use of \acrshort{glo:mcp} \cite{nsa2026}.

\par In public discussion these problems are usually grouped under the heading of ``AI security''. For part of them the label fits. Some attacks, such as hiding instructions inside the description of a tool, only work because a language model reads that text and acts on it. Many other reports, however, describe flaws that have been known for decades: a server that inserts user input into a shell command, a file tool that can be tricked into leaving its directory, a debugging service that listens on all network interfaces without a password. While reading the first advisories for this topic, I had the impression that the second kind was more common than the discussion suggested. If this impression is right, then the defences already exist and are well documented, and the real question is why they were skipped while the ecosystem was growing.

\par \textbf{Motivation and relevance.} For a software engineer the topic is practical. \acrshort{glo:mcp} servers are written by ordinary developers, very often as small open-source projects, and they carry the habits and shortcuts of any code written under time pressure. Meanwhile, companies and public institutions are planning to connect AI assistants to their internal systems. In Moldova this question became concrete in October 2026, when the e-government infrastructure suffered a data breach caused by an error in how access rights were checked \cite{tv82026}. The \acrshort{glo:mcp} security guidance warns about a related class of errors, in which an intermediary is trusted in place of the party that actually makes the request \cite{mcpsecbp}.

\par \textbf{Objective and research questions.} The objective of the article is to establish where \acrshort{glo:mcp} security failures originate. The work is organised around three questions:
\begin{itemize}
    \hitem \textbf{RQ1:} What share of the documented \acrshort{glo:mcp} weaknesses are classic software flaws, and what share depend on the behaviour of a language model?
    \hitem \textbf{RQ2:} At which stage (protocol design, implementation, integration, supply chain or deployment) were these weaknesses introduced?
    \hitem \textbf{RQ3:} Which signs of rushed engineering can be seen in the records, and what should teams that adopt \acrshort{glo:mcp} do about them?
\end{itemize}

\par \textbf{Methodology.} The study is a structured review of secondary data. I collected \nRecords{} documented weaknesses and incidents from a public vulnerability catalogue and from the original research reports, classified each of them using fixed rules, and compared the results with large measurement studies published in 2026. No real systems were scanned or tested.
